\begingroup
\raggedright
\subsection{Subset Sum, Knapsack and Max-Plus: Six Variants}
\label{sec:knapsack-six}
These are six useful parameter regimes from
\href{https://codeforces.com/blog/entry/98663}{errorgorn's tutorial},
with independently written implementations, derivations and examples.
Choose the method by its assumptions, not just by the word ``knapsack''.
All weights below are \emph{positive integers}. Subset sum asks for an
\emph{exact} target. The bounded, distinct-weight and window knapsack
implementations return the best value with weight \emph{at most} the
capacity; choosing nothing is allowed. NDPC exact-fill knapsack is a
separate subsection afterward.

\textbf{Notation and implementation conventions.}
$N$ is the input length, $C$ a target/capacity, $S=\sum_i w_i$,
$D$ a maximum weight, and $d$ a number of distinct weights.
These meanings are restated where needed.
The C++17 headers use namespace \texttt{knapsack\_ref}; include the
header corresponding to the method. Only the NDPC listing contains
\texttt{main}. Array sizes and integer index arithmetic must fit
\texttt{int} and available memory. Knapsack values, sums, differences
and products in these generic headers must fit \texttt{long long}.
For larger values, change the value type and its numeric limits together.

\begin{center}
\begin{tabular}{p{0.40\columnwidth}p{0.49\columnwidth}}
\textbf{Assumption / goal} & \textbf{Method / time} \\
\hline
All subset sums; total $S$ manageable & Coalesce equal weights, then bitsets; $O(N+S+(S/64+1)\sqrt S)$ \\
One subset-sum target; $w_i\le D$ & Balanced states; $O(ND)$ \\
Bounded copies per type & Residues and a monotone deque; $O(NC)$ \\
One concave max-plus operand & SMAWK; $O(n+m)$ \\
0/1 items; $d$ distinct weights & Group and convolve; $O(N\log N+dC)$ \\
Unlimited copies; $w_i\le D$, huge $C$ & Repeatedly double a short window; $O(N+D^2\log(C+2))$ \\
\end{tabular}
\end{center}

\subsubsection{1. All Subset Sums: Compress Repeated Weights}
\textbf{Problem.} Each of $N$ items may be taken once. Determine every
attainable sum from $0$ to $S=\sum_i w_i$. A boolean array costs $O(NS)$.
A bitset packs 64 states into a word: initially only bit zero is set,
and one item of weight $w$ performs
\[
 B\gets B\mathbin{\vert}(B\ll w).
\]
The shift uses the \emph{old} $B$, so the item is used at most once.
For a manually stored bitset, update words from high to low for the
same reason. A shift by an exact multiple of 64 needs a separate
case: shifting a 64-bit word by 64 is undefined.

\textbf{Reduce the number of shifts.} Let $c_w$ be the number of
remaining copies of weight $w$. Process weights in increasing order.
For odd $c_w$, retain one copy of $w$ and replace the other copies by
$(c_w-1)/2$ bundles of weight $2w$. For positive even $c_w$, retain two
copies and create $(c_w-2)/2$ bundles of weight $2w$.
Add these bundles to the count for $2w$, which will be processed later.
A bundle is one binary choice representing two copies, not an
unbounded item.

\textbf{Why all sums survive.} With one retained copy, selecting
$e\in\{0,1\}$ singles and $t$ pairs gives $e+2t$ original copies,
covering every count from zero through the original odd $c_w$.
With two retained copies, $e\in\{0,1,2\}$ covers both parities and
every count through the original even $c_w$. Conversely, every new
selection uses at most $c_w$ original copies. Thus the replacement
preserves exactly the attainable sums; applying it repeatedly does too.

\textbf{Why only $O(\sqrt S)$ shifts.} Total weight is conserved.
If $k$ distinct positive weights remain, their total is at least
$1+2+\cdots+k=k(k+1)/2\le S$. At most two copies of each survive,
so there are at most $2k=O(\sqrt S)$ bitset updates.
Counting/scanning takes $O(N+S)$; the shifts take
$O((S/64+1)\sqrt S)$. The implementation uses $O(S)$ integer counters
and $O(S/64+1)$ bitset words. This is a bound in terms of the
\emph{total sum}, not an arbitrary small target.

\textbf{Example.} Eight copies of weight two become two copies of two,
three copies of four; processing the latter leaves one four and one
eight. The final weights $[2,2,4,8]$ attain exactly
$0,2,4,\ldots,16$. Add an original weight three and one bitset shift
also gives the corresponding odd sums. Zero-weight items may be
ignored for feasibility; negative weights require another method.
\lstinputlisting[style=C++,showstringspaces=false,firstline=1,lastline=33]{../dp/SubsetSumSpeedups.hpp}
\begin{minipage}{\linewidth}
\begin{lstlisting}[style=C++,showstringspaces=false]
auto B = knapsack_ref::subset_sum_all({2, 2, 3});
int target = 7;
bool possible = (B[target / 64] >> (target % 64)) & 1;
// Check 0 <= target <= total sum before indexing.
// Attainable sums: 0, 2, 3, 4, 5, 7.
\end{lstlisting}
\end{minipage}
\par

\subsubsection{2. One Subset-Sum Target: Balanced States in O(ND)}
\textbf{Problem.} Each item may be used once, $1\le w_i\le D$, and we
only need to decide whether target $C$ is attainable. Here $D$ is
\emph{maximum item weight}, not the number of distinct weights.
The running time is $O(ND)$, independent of the numerical size of $C$.
Return true for $C=0$, and false when $C>S$.

\textbf{Start near the target.} Use zero-based indices. Take the
longest prefix $[0,k)$ whose sum $P$ is strictly smaller than $C$.
Because the next item has weight at most $D$, $C-D\le P<C$.
Any target solution can be reached from this prefix by removing some
prefix items and adding some suffix items.

There is an order of these operations that stays in the narrow band
$[C-D,C+D]$. Below $C$, add a required suffix item; above $C$, remove
a prefix item that is absent from the desired solution. The required
kind of operation always exists: otherwise the remaining operations
could not bring the current sum to $C$. Every step changes the total
by at most $D$. We may choose additions in increasing index order and
removals in decreasing index order; the same argument still applies
because any remaining item of the required kind is suitable.
Stop immediately if the total is $C$.

\textbf{Compress the choice of removed prefix.} After processing
suffix indices through $r$, store a boundary $l\in[0,k]$ for each
band total $s$. Indices $[0,l)$ of the initial prefix are still fixed
as selected. The last removed prefix index, if any, is $l$; future
removals may use only indices smaller than $l$. Initially $l=k$.
All already processed suffix choices and prefix indices at least $l$
are accounted for in the state's witness.

For the same $(s,r)$, retain only the \emph{largest} feasible $l$:
it permits every future removal allowed by a smaller boundary.
Future suffix choices depend only on $r$, not on the history.
Equivalently, one can define feasibility with $[0,l)$ fixed and the
processed interval after it flexible; feasibility for a larger
boundary implies feasibility for a smaller one. This is the dominance
that replaces a third DP dimension by one integer.

\textbf{Transitions and their order.}
\begin{enumerate}
\item Copy the previous suffix layer, representing skipping item $r$.
From each old total below $C$, also add $w_r$, keeping its boundary.
Only old-layer states are used, so item $r$ cannot be added twice.
\item Scan totals above $C$ from high to low. From boundary $l$,
removing fixed prefix item $j<l$ creates total $s-w_j$ and boundary
$j$. Descending totals ensure that if this new total is still above
$C$, its own removals are processed later in the same scan.
\item At total $s$, let $l_{\rm old}$ be the boundary before adding
this suffix item. Try only
$\max(0,l_{\rm old})\le j<l_{\rm new}$.
Indices smaller than $l_{\rm old}$ were already tried from this total
in an earlier layer, and their results remain available because
skipping later suffix items is permitted.
\end{enumerate}
For a fixed total, boundaries never decrease as suffix items are
processed. The newly tried ranges of $j$ therefore never overlap.
There are $2D+1$ totals and at most $k$ removals per total over the
\emph{whole run}, giving $O(ND)$ time. Two layers use $O(D)$ space.
The array index for total $s$ is $s-C+D$; index $D$ is the target.
A value $-1$ means unreachable, while boundary zero is reachable.

\textbf{Example.} Weights $[4,3,5]$, target eight: the initial prefix
has total seven. Add five to obtain twelve, then remove the prefix
item four to obtain eight. A DP that kept only totals at most eight
would discard the necessary intermediate total twelve.
\lstinputlisting[style=C++,showstringspaces=false,firstline=35,lastline=68]{../dp/SubsetSumSpeedups.hpp}
\begin{minipage}{\linewidth}
\begin{lstlisting}[style=C++,showstringspaces=false]
using knapsack_ref::subset_sum_small_weights;
bool yes = subset_sum_small_weights({4, 3, 5}, 8); // true
bool no  = subset_sum_small_weights({4, 3, 5}, 6); // false
\end{lstlisting}
\end{minipage}
\par
This is an exact deterministic algorithm. Randomly shuffling signed
items and retaining an experimentally chosen band is a different,
heuristic approach without this guarantee. The balancing technique
originates with Pisinger; see
\href{https://doi.org/10.1006/jagm.1999.1034}{Linear Time Algorithms
for Knapsack Problems with Bounded Weights (1999)}.

\subsubsection{3. Bounded Knapsack: Residues and a Monotone Deque}
\textbf{Problem.} Type $i$ has weight $w_i$, value $v_i$ and at most
$K_i$ copies; maximize value with weight at most $C$.
Here $N$ counts \emph{types}, not all expanded copies.
Let $F[c]$ be the best value with budget $c$ before this type, and
$G[c]$ the answer afterward. Initially $F[c]=0$ for every $c$, since
unused budget is permitted. The ordinary transition is
\[
 G[c]=\max_{0\le t\le K_i,\;tw_i\le c}
       \{F[c-tw_i]+t v_i\}.
\]
Trying every copy count costs too much when $K_i$ is large.

\textbf{Fix one residue.} Write $c=r+t w_i$, with $0\le r<w_i$.
If $s$ is the number of weight steps assigned to the old DP, then
$t-s$ is the number of new copies:
\[
 G[r+t w_i]=t v_i+
 \max_{\max(0,t-K_i)\le s\le t}
       \{F[r+s w_i]-s v_i\}.
\]
For increasing $t$, this is a sliding maximum over a window of
$K_i+1$ candidates. Store pairs $(s,F[r+s w_i]-s v_i)$ in a deque.
Before querying, remove expired indices $s<t-K_i$ from the front.
When inserting a new candidate, remove back entries with no larger
value: the newer candidate is at least as valuable and expires later.
The front is then the maximum. Insert $s=t$ before querying so taking
zero copies is included. Every index enters and leaves at most once.

\textbf{Why use an old layer?} Reading partially updated $G$ would
allow extra copies of this type. Copy $F$ before the residue loops.
All residues together contain $C+1$ states, so one type takes $O(C)$,
all types $O(NC)$, with $O(C)$ extra space.
Replace $K_i$ by $\min(K_i,\lfloor C/w_i\rfloor)$ first.
Types with $w_i>C$ cannot be taken.

\textbf{Example.} Suppose the old DP is all zero, a type has
$(w,v,K)=(3,5,2)$, and $C=8$. In residue two, budgets $2,5,8$ give
$t=0,1,2$. The candidates are $0,-5,-10$, and the resulting values
are $0,5,10$. At budget eleven, the old candidate $s=0$ would expire:
three copies are forbidden even though they would fit the budget.
\lstinputlisting[style=C++,showstringspaces=false]{../dp/BoundedKnapsack.hpp}
\begin{minipage}{\linewidth}
\begin{lstlisting}[style=C++,showstringspaces=false]
auto dp = knapsack_ref::bounded_knapsack(
    8, {{3, 5, 2}, {2, 4, 1}});
// dp[8] = 14: two weight-3 copies and one weight-2 copy.
\end{lstlisting}
\end{minipage}
\par

\textbf{Optional pruning before the DP.} For a fixed weight $w$,
no solution uses more than $\lfloor C/w\rfloor$ copies in total,
even if many types have that weight. Sort those types by decreasing
value and retain only the best that many copies, reducing the final
type's count if necessary. Do this using counts; do not expand huge
multiplicities. Replacing a retained low-value copy by a higher-value
copy of the same weight can only improve a solution, so the pruning
is safe. Across all weights there are at most
\[
 \sum_{w=1}^{C}\lfloor C/w\rfloor=O(C\log(C+1))
\]
retained copies, hence at most that many nonempty types. Combining
pruning with the deque DP gives
$O(N\log N+C\min(N,C\log(C+1)))$ time, including a simple sort.
The displayed function is the $O(NC)$ core; apply this pruning
separately when it reduces the number of types.

\subsubsection{4. Max-Plus Convolution with Concavity}
\textbf{Problem.} Given arrays $A$ of length $n$ and $B$ of length
$m$, compute
\[
 H[k]=\max_{\substack{0\le i<n\\0\le k-i<m}}
           (A[i]+B[k-i]),\quad 0\le k<n+m-1.
\]
This combines two independent DP choices whose sizes add to $k$.
The generic double loop costs $O(nm)$. Ordinary FFT convolution uses
addition and multiplication, so it does not directly compute this
maximum-and-addition operation.

\textbf{Use the correct curvature.} A sequence $B$ is discretely
\emph{concave} if $B[t+1]-B[t]$ is nonincreasing. For example,
$[0,9,15,16]$ has differences $[9,6,1]$ and is concave.
A convex sequence has nondecreasing differences.
The fast maximization methods here require concavity; negating both
operands gives the corresponding \emph{min-plus} method for convex
operands. A finite sequence of length one is both convex and concave.

\textbf{Case A: both operands are concave.} Start with
$H[0]=A[0]+B[0]$. Merge the two already sorted lists of differences
in decreasing order and append each difference to the running value.
For any $k$, the largest $k$ differences of their union maximize the
sum. Since each list is decreasing, the selected differences can be
chosen as a prefix of each list; such prefixes correspond to legal
indices $i$ and $k-i$. This proves the merge rule in $O(n+m)$ time.
For $A=[0,8,11]$, $B=[0,6,7]$, merge $[8,3]$ and $[6,1]$ to obtain
$[8,6,3,1]$, hence $H=[0,8,14,17,18]$.

\textbf{Case B: only $B$ is concave.} Build an implicit matrix with
row $k$ and column $i$:
\[
 M[k,i]=A[i]+B[k-i],\qquad i\le k<i+m.
\]
A row maximum is exactly $H[k]$. We evaluate entries on demand;
materializing the whole matrix would lose the speedup.
For columns $i<j$ where both entries exist, define the advantage of
$j$ over $i$ as
\[
 g(k)=A[j]-A[i]+B[k-j]-B[k-i].
\]
Let $\Delta B[t]=B[t]-B[t-1]$. Then
\[
 g(k+1)-g(k)=\Delta B[k+1-j]-\Delta B[k+1-i]\ge0,
\]
since the earlier difference of a concave sequence is no smaller.
Thus once a right column is at least as good as a left column,
it stays so in later rows. This property holds for every submatrix:
the matrix is \emph{totally monotone for row maxima}. With ties
resolved to the right, maximizing columns are nondecreasing.
No concavity assumption is made on $A$.

\textbf{Boundary comparisons matter.} Legal columns of row $k$ form
$[\max(0,k-m+1),\min(n-1,k)]$. Assigning all illegal cells the same
numeric minus infinity can destroy total monotonicity in submatrices
containing only illegal entries. Instead compare ordered columns
$i<j$ as follows:
\begin{enumerate}
\item If $j>k$, prefer $i$: $j$ is too far right.
\item Otherwise, if $k-i\ge m$, prefer $j$: $i$ is too far left.
\item Otherwise both entries are valid; compare their values,
preferring $j$ on equality.
\end{enumerate}
For two fixed columns this preference still switches only from left
to right as the row increases. Every full row has a legal column,
so its reported optimum is legal. The implementation accepts finite
values; its knapsack caller uses budget states initialized to zero,
so unreachable-state sentinels are unnecessary.

\textbf{How SMAWK finds all maxima in linear time.}
\begin{enumerate}
\item \emph{Reduce columns.} Process columns from left to right,
keeping a stack of candidates. Compare a new column against the
stack top at the row corresponding to the stack's current height.
If the new column wins, pop the old one. Otherwise retain it when
there is space. Total monotonicity makes a popped column unnecessary:
later rows prefer the new column, and earlier rows have their
candidates protected by earlier stack positions. Keep at most one
candidate per row. Each column is pushed/popped at most once.
\item \emph{Recurse on odd-positioned rows.} Solve rows at positions
$1,3,5,\ldots$ using the reduced columns. Their maximizing columns
are in nondecreasing order.
\item \emph{Interpolate even-positioned rows.} A row between two
solved rows searches only between their maximizing columns. At an
end, use the first/last candidate as the missing bound. These search
intervals overlap only at endpoints, so the total scan is linear.
\end{enumerate}
The first reduction costs $O(n+m)$. Recursion halves the row count
and has at most the previous row count in columns; remaining work
forms a geometric series. Time and memory are $O(n+m)$.
The following header includes both the SMAWK version and the shorter
two-concave-array merge version.
\lstinputlisting[style=C++,showstringspaces=false]{../dp/ConcaveMaxPlus.hpp}
\begin{minipage}{\linewidth}
\begin{lstlisting}[style=C++,showstringspaces=false]
using knapsack_ref::max_plus_concave;
auto H = max_plus_concave({0, 4, 10}, {0, 7, 10});
// B differences: 7, 3. A need not be concave.
// H = {0, 7, 11, 17, 20}.
\end{lstlisting}
\end{minipage}
\par
For comparison, the same monotone maximizing indices also allow a
simpler divide-and-conquer search: solve the middle row, restrict
left rows to columns no larger than its optimum, and right rows to
columns no smaller. That is generally logarithmically slower;
SMAWK is what gives the linear bound used in the next variant.

\subsubsection{5. 0/1 Knapsack with Few Distinct Weights}
\textbf{Problem.} There are $N$ individual items, each available once,
capacity $C$, but only $d$ distinct usable weights. Processing every
item separately costs $O(NC)$; group by weight to obtain
$O(N\log N+dC)$ including sorting.

\textbf{Summarize one weight group.} Suppose its weight is $w$ and
its values in descending order are $v_1\ge v_2\ge\cdots\ge v_K$.
If exactly $t$ copies from this group are taken, the best choice is
the $t$ largest values. Define
$P[0]=0$ and $P[t]=\sum_{j=1}^{t}v_j$.
Since $P[t]-P[t-1]=v_t$ is nonincreasing, $P$ is concave.
Keep only $\min(K,\lfloor C/w\rfloor)$ values: additional copies
cannot fit. Values may be negative because choosing zero copies
remains an option.

\textbf{Derive the convolution.} Let $F[c]$ be the previous best
value with budget $c$ and $G[c]$ the new answer. For residue
$r\in[0,w)$, form $A[s]=F[r+sw]$. Then
\[
 G[r+tw]=\max_{0\le j\le\min(t,K)}
                \{A[t-j]+P[j]\}=(A\star_{\max,+}P)[t].
\]
Thus call the one-concave-operand convolution from variant 4 and
retain only entries corresponding to budgets at most $C$.
Build every residue from an \emph{old} copy of $F$ so the group
cannot be reused. $F[c]=0$ initially for every budget, making all
$A$ entries finite and implementing ``weight at most'' consistently.

\textbf{Account for all work.} The lengths of all $A$ arrays for one
weight sum to $C+1$. Repeating $P$ in $w$ residue calls costs at most
$w(K+1)\le C+w=O(C)$ because $w\le C$.
The full convolution outputs are also $O(C)$ in total. Therefore
one weight group costs $O(C)$ even though it makes several calls.
The arbitrary old array $A$ usually is not concave; use SMAWK, not
the merge of two sorted difference lists. Memory is $O(N+C)$ for
the groups and DP, with $O(C)$ temporary convolution storage.

\textbf{Example.} Capacity eight; weight-two values $[9,6,1]$ give
$P=[0,9,15,16]$. After this group, the budget DP is
$[0,0,9,9,15,15,16,16,16]$.
Now process a weight-three item worth eight. At budget eight,
take the two best weight-two items plus the weight-three item:
weight seven, value $9+6+8=23$. Budget states correctly allow the
one unused unit.
\lstinputlisting[style=C++,showstringspaces=false]{../dp/DistinctWeightKnapsack.hpp}
\begin{minipage}{\linewidth}
\begin{lstlisting}[style=C++,showstringspaces=false]
auto dp = knapsack_ref::distinct_weight_knapsack(
    8, {{2, 9}, {2, 6}, {2, 1}, {3, 8}});
// dp[8] == 23. Each listed item can be taken only once.
\end{lstlisting}
\end{minipage}
\par

\subsubsection{6. Unbounded Knapsack with Small Weights, Huge Capacity}
\textbf{Problem.} Item types have $1\le w_i\le D$, unlimited copies,
and values $v_i$. Capacity $C$ can be far too large for an array of
length $C+1$. Return the optimum with total weight at most $C$ in
$O(N+D^2\log(C+2))$ time. No concavity is assumed here.
For duplicate weights, keep only the greatest value.

\textbf{Turn a budget into an exact sum safely.} Add a dummy item
of weight one and value zero, available without limit. Any unused
capacity can now be filled with dummy items, preserving value.
Conversely, removing dummy items leaves a feasible budget solution.
Let $F[s]$ be this padded exact-sum optimum. Every $s\ge0$ is
reachable, $F[0]=0$, and $F[C]$ is the desired answer.
The dummy is appropriate here because unused capacity is allowed;
it must not be added to the NDPC exact-fill problem.

\textbf{Balanced partition lemma.} Take any multiset of items, dummy
items included, and assign its items one at a time to the lighter
of two piles. If the difference in pile weights was at most $D$,
adding a weight at most $D$ to the lighter pile leaves a difference
at most $D$. Starting from two empty piles proves that a total $s$
can be partitioned into $a+b=s$ with $|a-b|\le D$.
Apply this to an optimal multiset. Each part's value is at most
$F[a]$ or $F[b]$, and combining two optimal parts is feasible.
Therefore
\[
 F[s]=\max_{\substack{a+b=s,\;a,b\ge0\\|a-b|\le D}}
                    \{F[a]+F[b]\}.
\]
Unlimited copies are essential: combining independently optimal
parts may repeat an item type arbitrarily often.

\textbf{Compute only a window.} Define a routine for center $T$
returning $F[s]$ for $\max(0,T-D)\le s\le T+D$.
If $T\le2D$, all these indices are at most $3D$; obtain them from
a normal unbounded DP precomputed through $3D$ in $O(D^2)$.
Otherwise recursively compute the window centered at
$h=\lfloor T/2\rfloor$. For any $s\in[T-D,T+D]$, the balanced
split obeys
\[
 \left\lceil\frac{s-D}{2}\right\rceil\le a,b\le
 \left\lfloor\frac{s+D}{2}\right\rfloor,
\]
so both indices lie in $[h-D,h+D]$, including when $T$ is odd.
Convolve the child's window with itself using the ordinary double
loop. If its first absolute index is $L$, entries $i,j$ produce
absolute total $s=2L+i+j$. Keep only totals in the parent's window.
Taking \emph{all} child pairs is safe: they are feasible solutions,
and the balanced partition lemma ensures an optimal pair is present.

\textbf{Why logarithmic in capacity.} The centers are
$C,\lfloor C/2\rfloor,\lfloor C/4\rfloor,\ldots$ until at most $2D$.
Each window has at most $2D+1$ entries, each double loop costs
$O(D^2)$, and there is only one recursive child at each level.
This gives the claimed time. The code releases the child at each
return, using $O(D+\log(C+2))$ extra memory rather than storing
all windows. It supports $D\ge1$, $C\le10^{18}$ and integer index
sizes large enough for $3D+1$.

\textbf{Example.} For types $(w,v)=(2,5),(3,8)$, capacity ten has
optimum 26 from $3+3+2+2$. For capacity $41$ and $D=3$, the parent
needs totals $[38,44]$ and its child, centered at 20, stores
$[17,23]$. Every total in the parent has a balanced split entirely
inside the child. This avoids computing all answers up to 41 and
works the same way when 41 is replaced by $10^{18}$.
\lstinputlisting[style=C++,showstringspaces=false]{../dp/UnboundedKnapsackWindow.hpp}
\begin{minipage}{\linewidth}
\begin{lstlisting}[style=C++,showstringspaces=false]
long long ans = knapsack_ref::unbounded_knapsack_window(
    10, 3, {{2, 5}, {3, 8}}); // 26
\end{lstlisting}
\end{minipage}
\par
The grouped-weight and short-window techniques are also developed in
\href{https://arxiv.org/abs/1802.06440}{Axiotis and Tzamos,
Capacitated Dynamic Programming: Faster Knapsack and Graph Algorithms}.
The implementations here use explicit integer window endpoints and
budget-state conventions throughout.

\subsection{NDPC N: Exact-Fill Knapsack with Many Huge Queries}
\label{sec:ndpc-knapsack}
\textbf{Problem and bounds.}
\href{https://atcoder.jp/contests/ndpc/tasks/ndpc2026_n}{Next DP
Contest, N -- Knapsack}: type $i$ has weight $i$ and value $v_i$,
for $1\le i\le n$; every type has unlimited copies. Each query asks
for maximum value at \emph{exactly} weight $W$.
Official bounds are $n\le4000$, $q\le2\cdot10^5$,
$0\le v_i\le10^9$, $1\le W\le10^9$.
A weight-one type exists, so every nonnegative exact weight is
reachable. The listing below implements the supplied algorithm,
with explicit-width density comparisons and a guarded unreachable
state. Under these bounds, 64-bit values suffice because an answer
is at most $W\max_i v_i\le10^{18}$. Outside these bounds, recheck
both arithmetic and the $O(n^2)$ memory requirement.

\subsubsection*{Step 1: Sort Types by Value per Unit Weight}
An item $a$ precedes $b$ when $v_a/a>v_b/b$. Compare
$v_a b$ and $v_b a$ using integer arithmetic; do not divide or compare
floating-point ratios. Break equal-density ties by smaller original
index. Let \texttt{order} contain the original indices in this order,
and let $p=\texttt{order[0]}$. The weight of the best-density type
is $p$, not its position zero in the sorted list.
Ties cause no difficulty: replacements need only preserve value,
not strictly increase it.

\subsubsection*{Step 2: Reduce Every Large Query to at Most n Squared}
Let $L=n^2$. For every exact weight $W\ge L$, there exists an optimum
containing type $p$. To prove this, an optimum uses at least $n$
items because every item weighs at most $n$. Among any $p$ of those
items, the $p+1$ prefix sums have two equal residues modulo $p$.
The nonempty block between them weighs a positive multiple of $p$.
Replace that block by copies of type $p$. Its weight stays equal,
and its value does not decrease because $p$ has maximum density.
Thus an optimum containing $p$ exists, even if a different optimum
originally omitted it.

Writing $F[W]$ for the exact optimum, this proves
$F[W]=F[W-p]+v_p$ whenever $W\ge L$.
The forward inequality follows by removing $p$ from that optimum;
the reverse follows by adding $p$ to an optimum for $W-p$.
Repeatedly apply the equality until the remaining weight is at most
$L$. For $W>L$, the number of removed copies is
\[
 t=\left\lceil\frac{W-L}{p}\right\rceil,\qquad R=W-tp.
\]
Then $L-p<R\le L$, and the answer is $F[R]+t v_p$.
For $W\le L$, set $t=0,R=W$. This explains both the ceiling and
why its argument is $W-L$, rather than $W$ itself. The large part
can use the best-density item, but greedily filling \emph{all} of a
small query with it is not valid.

\subsubsection*{Step 3: Avoid n Transitions at Every Small Weight}
The ordinary exact-fill DP is
\[
 F[0]=0,\qquad F[w]=\max_{1\le i\le\min(n,w)}
                            \{F[w-i]+v_i\}.
\]
Computing it through $L$ costs $O(n^3)$. A density-rank lemma lets
us restrict the candidate final item. For exact weight $w$, some
optimum uses a type among the first
\[
 K(w)=\min\left(n,\left\lceil\frac{2n^2}{w}\right\rceil\right)
\]
entries of \texttt{order}. Remove that item from this optimum; its
remainder is an optimum of smaller exact weight, by the usual
optimal-substructure argument. Consequently it suffices to try these
$K(w)$ types in the recurrence, ignoring any whose weight exceeds
$w$. Increasing $w$ makes every needed predecessor already available.
This is a statement about \emph{at least one item} in an optimum;
it does not say all its items have high density.

\textbf{Why the rank lemma holds.} The nontrivial number-theoretic
ingredient is the following classical Frobenius bound:
for $h\ge2$ distinct positive integers with gcd one, maximum $M$,
and second largest $M_2$, the largest unrepresentable integer is
at most $2\lfloor M/h\rfloor M_2-M$ (Erd\H{o}s--Graham).
In particular it is at most $2M^2/h-M$.
Here is how that bound yields the exact statement used by the code.

Fix a rank $k<n$ and assume $w\ge2n^2/k$. Take any optimum omitting
the first $k$ types; if none exists the claim already holds.
Let $q$ be its highest-density used type, and let
$g=\gcd(w_1,\ldots,w_k,w_q)$, where these $w_j$ denote the weights
of the sorted types. Divide those $k+1$ distinct weights by $g$,
obtaining a gcd-one set $A$ with maximum $M\le n/g$.
For any chosen top-$k$ weight $a\in A$, the Frobenius bound gives
\[
 F_A+a\le\frac{2M^2}{k+1}<\frac{2M^2}{k}.
\]
Thus any integer $x\ge2M^2/k$ can be represented using $A$ while
including a top-$k$ weight: first reserve $a$, then represent $x-a$.
Here $F_A$ means the Frobenius number, not the knapsack DP.

We need a sufficiently large weight divisible by $g$ to replace.
As long as at least $g$ items remain in the optimum, equal residues
of $g+1$ prefix sums identify a nonempty block divisible by $g$.
Remove such blocks repeatedly and let their combined weight be $B$.
At most $g-1$ items remain, so
\[
 B\ge w-(g-1)n\ge\frac{2n^2}{gk}.
\]
For the last inequality, subtract the right side from
$2n^2/k-(g-1)n$ to get
$(g-1)n(2n/(gk)-1)\ge0$.
Indeed, $k+1$ distinct positive multiples of $g$ are at most $n$,
so $g(k+1)\le n$. Therefore
$B/g\ge2n^2/(g^2k)\ge2M^2/k$, and $B/g$ is an integer.
Represent it by $A$ while including a top-$k$ weight, then multiply
back by $g$. This replaces exactly weight $B$ using only top-$k$
types and type $q$. Every replacement type has density at least
that of $q$, which was highest among the removed items, so value
cannot decrease. We have constructed an optimum using a top-$k$
type, proving the claim. If $k=n$, it is immediate.
Taking $k=K(w)$ gives the restricted recurrence above.

This proof uses the classical Frobenius theorem as an external
lemma; the algorithmic reduction and its inequalities are shown
explicitly. See
\href{https://arxiv.org/abs/2110.02503}{Chan and He, More on
Change-Making and Related Problems}, Sections 4--5.
The distinction between exact weight and an at-most capacity is
necessary: the same factor two must not be applied blindly to
other knapsack definitions.

\subsubsection*{Step 4: Complexity, Code and a Query Walkthrough}
The transition count is
\[
 \sum_{w=1}^{n^2}\min\left(n,
              \left\lceil\frac{2n^2}{w}\right\rceil\right)
 =O(n^2\log(n+1)).
\]
For small $w$ use the cap $n$; for the rest use the harmonic-series
bound $\sum_{w=a}^{b}1/w=O(\log(b/a)+1)$.
Sorting costs $O(n\log n)$, queries $O(1)$ each, and memory is
$O(n^2)$. At $n=4000$, the 64-bit DP alone occupies about 128 MB
(decimal), so this is not a small-memory method.

\textbf{Example.} Let $n=3$ and $[v_1,v_2,v_3]=[2,5,7]$.
Densities are $2,2.5,7/3$, hence \texttt{order} is $[2,3,1]$,
$p=2$ and $L=9$. For query $W=20$,
$t=\lceil(20-9)/2\rceil=6$ and $R=8$.
The precomputed optimum is $F[8]=20$, giving
$20+6\cdot5=50$. For the small query $W=3$, the best answer is
seven; taking one best-density item and leaving one unit empty
would violate exact fill. The DP supplies the correct remainder.
\newpage
\lstinputlisting[style=C++,showstringspaces=false]{../dp/NDPC_Knapsack.cpp}
\begin{minipage}{\linewidth}
\begin{lstlisting}[style=C++,showstringspaces=false]
// Input:             Output:
// 3 4                7
// 2 5 7              20
// 3                  22
// 8                  50
// 9
// 20
\end{lstlisting}
\end{minipage}
\par
\par
\endgroup
